% Week 3 Session B — Test-Driven Development practice
% All lab reports in this series include an abstract.
%
% Build: pdflatex Week3_SessionB_Report.tex from Lab-reports/ (run twice)
% Screenshots (same folder):
%   week3_session_b_opencode_red.png, week3_session_b_opencode_green.png,
%   week3_session_b_opencode_refactor.png
%   week3_session_b_terminal_red.png, week3_session_b_terminal_green.png,
%   week3_session_b_terminal_refactor.png
% Code copies: week3_session_b_files/
%
\documentclass[11pt,a4paper]{article}

\usepackage[margin=2.5cm]{geometry}
\usepackage{graphicx}
\newcommand{\opencodefig}[1]{%
  \includegraphics[width=\linewidth,height=0.65\textheight,keepaspectratio]{#1}%
}
\newcommand{\terminalfig}[1]{%
  \includegraphics[width=0.92\linewidth,height=0.38\textheight,keepaspectratio]{#1}%
}
\usepackage{booktabs}
\usepackage{xurl}
\usepackage{hyperref}
\usepackage{parskip}
\usepackage{enumitem}
\usepackage{xcolor}
\usepackage{fvextra}
\usepackage[most]{tcolorbox}

\newcommand{\studentname}{Mahmudul Hasan}
\newcommand{\studentid}{4125999049}
\newcommand{\reportdate}{May 21, 2026}

\makeatletter
\g@addto@macro\UrlBreaks{\do\ }
\makeatother

\newcommand{\LabCodeRoot}{week3_session_b_files}

\newcommand{\filebox}[2]{%
\begin{tcolorbox}[
  enhanced,
  colback=gray!6,
  colframe=black!55,
  title={#1},
  fonttitle=\bfseries\ttfamily\footnotesize,
  arc=1.5mm,
  boxrule=0.7pt,
  breakable,
  width=\linewidth,
  before skip=0.8em,
  after skip=0.8em,
]
\VerbatimInput[fontsize=\footnotesize,breaklines=true,breakanywhere=true]{\LabCodeRoot/#2}
\end{tcolorbox}%
}

\title{Lab Report: Week 3 Session B\\Test-Driven Development Practice}
\author{\studentname \quad (Student ID: \studentid)}
\date{\reportdate}

\begin{document}
\maketitle

\begin{abstract}
This report documents Week~3 Session~B on \textbf{test-driven development (TDD)} with OpenCode Desktop: Red (failing tests before implementation), Green (minimum code to pass), and Refactor (improve structure while tests stay green). The SCS-CN function \texttt{calculate\_runoff(P, CN)} was developed in \texttt{week3\_session\_b/} separately from the Session~A scaffold. Red phase produced \texttt{ModuleNotFoundError}; Green and Refactor each yielded \textbf{7 passed} tests. Sanity checks $Q(50,80)\approx 13.80$~mm and $Q(80,85)\approx 43.55$~mm match earlier course labs. Evidence: six screenshots (OpenCode and terminal per phase) and full source in Appendix~\ref{app:files}--\ref{app:promptlog}.
\end{abstract}

\section{Introduction}
TDD treats tests as a \textbf{contract} for correct behavior: write tests first (Red), implement until they pass (Green), then refactor without changing outputs. This reduces wrong formulas and plausible-but-false AI outputs. Session~B applies TDD to the same SCS-CN physics as Session~A, but in a fresh folder \texttt{week3\_session\_b/} with tests written \emph{before} \texttt{src/runoff.py} existed.

\section{Objectives}
\begin{itemize}[noitemsep]
  \item Apply the Red $\rightarrow$ Green $\rightarrow$ Refactor cycle with AI assistance.
  \item Write comprehensive pytest tests before implementation.
  \item Use tests to validate AI-generated hydrology code.
  \item Document the process in \texttt{prompt\_log.md}.
\end{itemize}

\section{Methods}

\subsection{Conventions (all phases)}
$S = 25400/\mathrm{CN} - 254$~mm, $I_a = 0.2S$, $Q=0$ if $P \le I_a$, else $Q=(P-I_a)^2/(P-I_a+S)$, with $0 \le Q \le P$. Invalid $P<0$ or $\mathrm{CN}\notin[1,100]$ raise \texttt{ValueError}. Python~3.8 compatible hints (\texttt{Union}, not \texttt{X | Y}).

\subsection{Exercise 1: Red phase}
Prompt \texttt{prompt\_red.txt} asked OpenCode for \texttt{tests/test\_runoff.py} only---seven cases: normal ($P=50$, $\mathrm{CN}=80$), $P=0$, $P<I_a$, $\mathrm{CN}=100$, negative $P$, invalid $\mathrm{CN}$, and a $7\times7$ grid asserting $Q \le P$. No \texttt{src/runoff.py} was created.

\subsection{Exercise 2: Green phase}
Prompt \texttt{prompt\_green.txt} included the full test file; OpenCode implemented \texttt{src/runoff.py} to satisfy all tests.

\subsection{Exercise 3: Refactor phase}
Prompt \texttt{prompt\_refactor.txt} requested readability and clearer errors without changing numerical results; \texttt{pytest} was run after refactoring.

\section{Results}

\subsection{Submitted file inventory}
\begin{table}[htbp]
  \centering
  \caption{Files in \texttt{week3\_session\_b/}.}
  \label{tab:files}
  \small
  \begin{tabular}{@{}ll@{}}
    \toprule
    Path & Role \\
    \midrule
    \texttt{tests/test\_runoff.py} & Seven pytest cases (written in Red) \\
    \texttt{src/runoff.py} & SCS-CN implementation (Green, refactored) \\
    \texttt{prompt\_red.txt}, \texttt{prompt\_green.txt}, \texttt{prompt\_refactor.txt} & AI prompts \\
    \texttt{prompt\_log.md} & TDD documentation \\
    \texttt{requirements.txt} & pytest dependency \\
    \bottomrule
  \end{tabular}
\end{table}

\subsection{Red phase}
Terminal and OpenCode confirm collection failure: \texttt{ModuleNotFoundError: No module named 'src.runoff'} (Figures~\ref{fig:red-oc}--\ref{fig:red-term}).

\begin{figure}[htbp]
  \centering
  \opencodefig{week3_session_b_opencode_red.png}
  \caption{OpenCode Red: generated \texttt{test\_runoff.py}; confirmed import error.}
  \label{fig:red-oc}
\end{figure}

\begin{figure}[htbp]
  \centering
  \terminalfig{week3_session_b_terminal_red.png}
  \caption{Terminal Red: \texttt{pytest -v} --- \texttt{ModuleNotFoundError} (expected).}
  \label{fig:red-term}
\end{figure}

\subsection{Green phase}
All seven tests passed on first implementation (Figures~\ref{fig:green-oc}--\ref{fig:green-term}). Hand-checked: $Q(50,80)\approx 13.80$~mm; $Q(80,85)\approx 43.55$~mm.

\begin{figure}[htbp]
  \centering
  \opencodefig{week3_session_b_opencode_green.png}
  \caption{OpenCode Green: \texttt{src/runoff.py} created; 7 tests pass.}
  \label{fig:green-oc}
\end{figure}

\begin{figure}[htbp]
  \centering
  \terminalfig{week3_session_b_terminal_green.png}
  \caption{Terminal Green: \texttt{7 passed}.}
  \label{fig:green-term}
\end{figure}

\subsection{Refactor phase}
Refactoring introduced named constants, validation helpers, and \texttt{\_compute\_runoff}; public \texttt{calculate\_runoff} became a short validate$\rightarrow$compute$\rightarrow$clamp pipeline. Tests remained green (Figures~\ref{fig:ref-oc}--\ref{fig:ref-term}).

\begin{table}[htbp]
  \centering
  \caption{Refactoring changes (behavior unchanged).}
  \label{tab:refactor}
  \small
  \begin{tabular}{@{}p{0.32\linewidth}p{0.58\linewidth}@{}}
    \toprule
    Change & Rationale \\
    \midrule
    \texttt{\_RETENTION\_FACTOR}, etc. & Remove magic numbers \\
    \texttt{\_validate\_positive\_rainfall}, \texttt{\_validate\_curve\_number} & Separate validation; clearer errors \\
    \texttt{\_compute\_runoff} & Pure math on validated floats \\
    \texttt{effective\_rainfall} & Clearer than repeated $P-I_a$ \\
    Short \texttt{calculate\_runoff} & Readable public API \\
    \bottomrule
  \end{tabular}
\end{table}

\begin{figure}[htbp]
  \centering
  \opencodefig{week3_session_b_opencode_refactor.png}
  \caption{OpenCode Refactor: structural improvements; tests still pass.}
  \label{fig:ref-oc}
\end{figure}

\begin{figure}[htbp]
  \centering
  \terminalfig{week3_session_b_terminal_refactor.png}
  \caption{Terminal Refactor: \texttt{7 passed} after refactor.}
  \label{fig:ref-term}
\end{figure}

Table~\ref{tab:tdd} summarizes the TDD cycle.

\begin{table}[htbp]
  \centering
  \caption{TDD cycle outcomes.}
  \label{tab:tdd}
  \begin{tabular}{@{}lll@{}}
    \toprule
    Phase & \texttt{src/runoff.py} & \texttt{pytest} \\
    \midrule
    Red & Absent & Import error \\
    Green & Present (monolithic) & 7 passed \\
    Refactor & Helpers + constants & 7 passed \\
    \bottomrule
  \end{tabular}
\end{table}

\section{Analysis and validation}

\textbf{Tests as contract:} The normal-case expected value ($\approx 13.80$~mm) and impervious case ($Q=P$ when $\mathrm{CN}=100$) encode physics; a wrong $S$ formula or missing $P \le I_a$ branch would fail immediately.

\textbf{AI role:} OpenCode generated correct tests and implementation; refactor improved maintainability without changing outputs. Local \texttt{pytest} was still run to verify agent claims.

\textbf{Session A vs.\ B:} Session~A scaffolded a full project with tests after code; Session~B enforced tests-first TDD in an empty implementation folder.

\textbf{Hallucination check:} Pre-written tests constrained the Green-phase formula; Red prevented shipping untested code.

\section{Discussion}
\begin{enumerate}[noitemsep]
  \item Red phase failure was intentional and gradable evidence, not an environment bug.
  \item Green passed on the first try because tests fully specified behavior.
  \item Refactor improved error messages (e.g.\ ``got \{CN\}'') useful for debugging bad inputs.
  \item Next time: add one extra hand-calculated case in tests when trusting a new AI model.
\end{enumerate}

\section{Problems encountered}
Using a virtual environment from \texttt{week3\_session\_a} for \texttt{week3\_session\_b} worked but is optional; system \texttt{python3} and \texttt{pytest} were sufficient. Initial confusion that Red ``errors'' needed fixing was resolved by understanding TDD phase goals.

\section{Lessons learned}
TDD with AI still requires human verification: run \texttt{pytest} locally, keep tests unchanged during Refactor unless physics intentionally changes, and document each phase in \texttt{prompt\_log.md} with screenshots.

\section{Conclusion}
Week~3 Session~B objectives were met: Red, Green, and Refactor with SCS-CN runoff, seven passing tests throughout Green and Refactor, and documented prompts and outcomes. Submitted artifacts: \texttt{week3\_session\_b/}, \texttt{prompt\_log.md}, and phase screenshots.

\appendix
\section{Submitted project files}
\label{app:files}

\subsection{\texttt{tests/test\_runoff.py} (Red phase)}
\label{app:tests}
\filebox{\texttt{tests/test\_runoff.py} (full file)}{tests/test_runoff.py}

\subsection{\texttt{src/runoff.py} (Green + Refactor)}
\label{app:runoff}
\filebox{\texttt{src/runoff.py} (full file)}{src/runoff.py}

\subsection{Prompts}
\label{app:prompts}

\filebox{\texttt{prompt\_red.txt}}{prompt_red.txt}

\filebox{\texttt{prompt\_green.txt}}{prompt_green.txt}

\filebox{\texttt{prompt\_refactor.txt}}{prompt_refactor.txt}

\filebox{\texttt{requirements.txt}}{requirements.txt}

\section{Submitted prompt log: \texttt{prompt\_log.md}}
\label{app:promptlog}
\filebox{\texttt{prompt\_log.md} (full file)}{prompt_log.md}

\end{document}
